% GENERATED FILE. Edit canonical lessons, then run npm run generate:books.
% canonical-source-hash: fnv1a64:358566807a983f74
% canonical-lessons: BN-C21-ei, BN-C21-oi, BN-C21-dem-order, BN-C21-eta, BN-R21-pointing

\chapter{Pointing at Things}
\label{ch:bn-pointing-at-things}
\hlchaptermodality{32}

\begin{chapteropening}
\textbf{By the end of this chapter:} \emph{I can point at something near me and something across the room, put a pointer and a classifier on the same phrase, and ask what is this?}

The near and far heads arrive together, stack with the classifier at the opposite end of the phrase, and stand alone as \bn{এটা} — which with the \bn{কী} of chapter three makes the commonest beginner question in any language sayable for the first time.
\end{chapteropening}

\section[ei]{\bn{এই} (ei) — this}

\label{lesson:BN-C21-ei}



\begin{quote}
Say \textbf{the shirt}, then \textbf{a shirt}, then \textbf{another shirt}.

\begin{itemize}
\raggedright
  \item \emph{Recall:} two chapters back, the word that opens a clause about a moment — \emph{jôkhon}
\end{itemize}
\end{quote}

\subsection*{You'll want to know: \bn{এই}}
\begin{quote}
\textbf{\bn{এই} \bn{জামা}} — \emph{ei jāmā} — \textbf{this shirt.}
\end{quote}

\begin{quote}
\textbf{\bn{এই} \bn{চা}} — \emph{ei chā} — \textbf{this tea.}
\end{quote}

Thirty-one chapters, and nothing in the room could be pointed at. \textbf{\bn{টা}} says \emph{the one we both mean}; it does not say \emph{this one, here, the one I am looking at}. That is a different job and it needs a different word.

\textbf{\bn{এই}} goes \textbf{in front} of the noun — the opposite end from the classifier, and worth noticing, because those two are about to appear in the same phrase.

\begin{grammarlens}[title={Grammar Lens: the near head}]
Bengali builds a family of pointing words on one frame with three heads, and you met the frame in chapter thirty:

\noindent
\begin{tabularx}{\linewidth}{@{}>{\raggedright\arraybackslash}X>{\raggedright\arraybackslash}X>{\raggedright\arraybackslash}X@{}}
\toprule
\textbf{head} & \textbf{means} & \textbf{you have seen it in} \\
\midrule
\textbf{\bn{এ}-} & near — \emph{this} & \textbf{\bn{এই}} \\
\textbf{\bn{ও}-} & far — \emph{that} & the next lesson \\
\textbf{\bn{ক}-} & asking — \emph{which, where} & \textbf{\bn{কেমন}}, \textbf{\bn{কখন}} \\
\bottomrule
\end{tabularx}

So \textbf{\bn{এই}} is not a word to be memorised on its own. It is one square of a grid, and the rest of the grid arrives over the next six lessons.
\end{grammarlens}

\subsection*{Your turn}
\begin{itemize}
\raggedright
  \item \emph{Say it:} \emph{ei jāmā} — this shirt
  \item \emph{Say it:} \emph{ei chā} — this tea
  \item \emph{Say it:} the contrast — \emph{ekṭā jāmā}, then \emph{ei jāmā}
  \item \emph{Write it:} \bn{এই}
\end{itemize}

\begin{tcolorbox}[breakable,colback=teal!4,colframe=teal!35!black,title={Before you move on}]
Which end of the noun does \textbf{\bn{এই}} go on? (\textbf{The front} — unlike the classifier.) What does it say that \textbf{\bn{টা}} does not? (\textbf{Here, near me} — it points.) And which head of the pointing frame is it? (\textbf{\bn{এ}-}, the near one.)

Source: \href{https://www.unicode.org/charts/PDF/U0980.pdf}{Unicode Bengali chart}.
\end{tcolorbox}
\section[oi]{\bn{ওই} (oi) — that}

\label{lesson:BN-C21-oi}



\begin{quote}
Say \textbf{this shirt}. Then write \textbf{\bn{ও}}.
\end{quote}

\subsection*{You'll want to know: \bn{ওই}}
\begin{quote}
\textbf{\bn{ওই} \bn{জামা}} — \emph{oi jāmā} — \textbf{that shirt.}
\end{quote}

\begin{quote}
\textbf{\bn{এই} \bn{লাল}, \bn{ওই} \bn{কালো}} — \emph{ei lāl, oi kālo} — \textbf{this one red, that one black.}
\end{quote}

Swap the head and you have crossed the room. \textbf{\bn{এ}-} is here; \textbf{\bn{ও}-} is there. That is the whole system: Bengali has \textbf{two} distances, not three. Spanish insists on \emph{este}, \emph{ese} and \emph{aquel}; Bengali stops at two and lets context do the rest.

\begin{grammarlens}[title={Grammar Lens: one letter, two jobs}]
You have met \textbf{\bn{ও}} before, in chapter twenty-eight, as the joiner meaning \textbf{also}. It is the same letter and it is not an accident: a word that means \emph{that one over there} and a word that means \emph{and that too} are the same gesture, pointing away from the speaker.

Told apart by where it sits:

\begin{itemize}
\raggedright
  \item \textbf{\bn{ও}} alone, between two things — \emph{also, and}
  \item \textbf{\bn{ওই}} in front of a noun — \emph{that}
\end{itemize}

You will meet a third job for it two chapters from now, when it turns into a person — and it is still the same gesture.
\end{grammarlens}

\subsection*{Your turn}
\begin{itemize}
\raggedright
  \item \emph{Say it:} \emph{oi jāmā} — that shirt
  \item \emph{Say it:} the pair, pointing twice — \emph{ei jāmā}, \emph{oi jāmā}
  \item \emph{Say it:} \emph{ei lāl, oi kālo}
  \item \emph{Write it:} \bn{এই} \bn{ওই}
\end{itemize}

\begin{tcolorbox}[breakable,colback=teal!4,colframe=teal!35!black,title={Before you move on}]
How many distances does Bengali point to? (\textbf{Two} — here and there.) What is the only difference between \textbf{\bn{এই}} and \textbf{\bn{ওই}}? (\textbf{The head} — near against far.) And what else does \textbf{\bn{ও}} do, on its own? (\textbf{Also, and}.)

Source: \href{https://www.unicode.org/charts/PDF/U0980.pdf}{Unicode Bengali chart}.
\end{tcolorbox}
\section[ei jāmāṭā]{\bn{এই} \bn{জামাটা} (ei jāmāṭā) — this shirt, the one here}

\label{lesson:BN-C21-dem-order}



\begin{quote}
Say \textbf{this shirt}, then \textbf{that shirt}, then \textbf{the shirt}.
\end{quote}

\subsection*{You'll want to know: both ends at once}
\begin{quote}
\textbf{\bn{এই} \bn{জামাটা}} — \emph{ei jāmāṭā} — \textbf{this shirt} (the one here, and we both know which).
\end{quote}

\begin{quote}
\textbf{\bn{ওই} \bn{কাপড়টা}} — \emph{oi kāpôṛṭā} — \textbf{that cloth.}
\end{quote}

English cannot do this. \emph{The this shirt} is not English, because \emph{the} and \emph{this} are fighting over one slot in front of the noun and only one of them can win. In Bengali they are nowhere near each other:

\textbf{\bn{এই}} + \textbf{\bn{জামা}} + \textbf{\bn{টা}} · pointer · noun · classifier

Front and back. No competition, so no choice to make — and a Bengali speaker uses both far more often than an English speaker would expect, because leaving the classifier off sounds unfinished.

\begin{grammarlens}[title={Grammar Lens: the fixed order}]
Once a number and an adjective join in, the order still does not move:

\noindent
\begin{tabularx}{\linewidth}{@{}>{\raggedright\arraybackslash}X>{\raggedright\arraybackslash}X>{\raggedright\arraybackslash}X@{}}
\toprule
\textbf{slot} & \textbf{what goes there} & \textbf{example} \\
\midrule
first & pointer & \textbf{\bn{এই}} \\
second & number & \textbf{\bn{একটা}} \\
third & adjective & \textbf{\bn{লাল}} \\
last & noun + classifier & \textbf{\bn{জামাটা}} \\
\bottomrule
\end{tabularx}

Everything before the noun is in a fixed queue, and the only thing that ever attaches to the noun itself is the classifier at the back. Bengali will keep adding suffixes at that back end for the rest of this book, always one at a time, always in order.
\end{grammarlens}

\subsection*{Your turn}
Say these aloud:

\begin{itemize}
\raggedright
  \item \emph{ei jāmāṭā}
  \item \emph{oi kāpôṛṭā}
  \item both ends, then the front alone — \emph{ei jāmāṭā}, then \emph{ei jāmā}
  \item the long queue — \emph{ei lāl jāmāṭā}
\end{itemize}

\begin{tcolorbox}[breakable,colback=teal!4,colframe=teal!35!black,title={Before you move on}]
Can a pointer and a classifier appear in the same phrase? (\textbf{Yes} — they sit at opposite ends.) Why can English not do the same? (\textbf{Its \emph{the} and its \emph{this} want the same slot}.) And which of the two attaches to the noun itself? (\textbf{The classifier}, at the back.)

Source: \href{https://www.unicode.org/charts/PDF/U0980.pdf}{Unicode Bengali chart}.
\end{tcolorbox}
\section[eṭā]{\bn{এটা} (eṭā) — this one}

\label{lesson:BN-C21-eta}



\begin{quote}
Say \textbf{this shirt}, with both ends. Then say the question from chapter two that asks somebody their name.
\end{quote}

\subsection*{You'll want to know: \bn{এটা}}
\begin{quote}
\textbf{\bn{এটা} \bn{কী}?} — \emph{eṭā ki?} — \textbf{what is this?}
\end{quote}

\begin{quote}
\textbf{\bn{এটা} \bn{জল}} — \emph{eṭā jôl} — \textbf{this is water.}
\end{quote}

Take the noun out of \textbf{\bn{এই} \bn{জামাটা}} and keep the two ends. What is left — \textbf{\bn{এ}} and \textbf{\bn{টা}} — is \textbf{\bn{এটা}}, and it means \emph{this thing}, standing on its own as the subject of a sentence.

Chapter two taught \textbf{\bn{কী}} as the question word inside \emph{what is your name}, and then gave you nothing else to ask it about. \textbf{\bn{এটা}} is the missing half. Between them they make the first question a learner of any language actually uses, and it has been out of reach for thirty-one chapters.

The far one works the same way: \textbf{\bn{ওটা}} — \emph{oṭā} — \textbf{that one}.

\begin{grammarlens}[title={Grammar Lens: no verb in the middle}]
Notice what is \textbf{not} in \textbf{\bn{এটা} \bn{জল}}. There is no \emph{is}. Bengali joins a subject to what it is by putting the two side by side and stopping, exactly as chapter twenty-seven showed when it taught you to deny such a sentence with \textbf{\bn{নয়}}:

\begin{itemize}
\raggedright
  \item \textbf{\bn{এটা} \bn{জল}} — this is water
  \item \textbf{\bn{এটা} \bn{জল} \bn{নয়}} — this is not water
\end{itemize}

So the sentence and its denial both arrive at once, because the denier was already taught and had nothing to deny.
\end{grammarlens}

\subsection*{Your turn}
Say these aloud:

\begin{itemize}
\raggedright
  \item \emph{eṭā ki?}
  \item \emph{eṭā jôl}, then \emph{eṭā chā}
  \item the denial — \emph{eṭā jôl nôy}
  \item the far one — \emph{oṭā ki?}
  \item with a noun and without — \emph{ei jāmāṭā}, then \emph{eṭā}
\end{itemize}

\begin{tcolorbox}[breakable,colback=teal!4,colframe=teal!35!black,title={Before you move on}]
What two pieces is \textbf{\bn{এটা}} made of? (\textbf{\bn{এ}-}, the near head, and the classifier.) Is there a word for \textbf{is} in \textbf{\bn{এটা} \bn{জল}}? (\textbf{No} — the two halves sit side by side.) And which word denies it? (\textbf{\bn{নয়}}.)

Source: \href{https://www.unicode.org/charts/PDF/U0980.pdf}{Unicode Bengali chart}.
\end{tcolorbox}
\section[ei · oi · eṭā]{Pointing at things}

\label{lesson:BN-R21-pointing}



\begin{quote}
How many distances does Bengali point to, and what is the only piece that changes between them? Then say \textbf{this one} and \textbf{that one} with no noun at all.
\end{quote}

\subsection*{Your turn}
\begin{itemize}
\raggedright
  \item \emph{Say it:} \emph{ei jāmāṭā}, then \emph{oi jāmāṭā}
  \item \emph{Say it:} \emph{eṭā ki?} — and answer it
  \item \emph{Say it:} the denial — \emph{eṭā chā nôy}
  \item \emph{Recall:} from the chapter before, the ending that means \emph{the} — \emph{-ṭā}
\end{itemize}

\subsection*{Your turn: the distant band — the room, pointed at}
Two chapters ago none of this could be said. Take the whole book's nouns and point at them, near and far, cold.

Say these aloud:

\begin{itemize}
\raggedright
  \item the drinks, near then far — \emph{ei chāṭā}, \emph{oi chāṭā}; \emph{ei jôlṭā}, \emph{oi dudhṭā}
  \item the wardrobe with a colour in the queue — \emph{ei lāl jāmāṭā}, \emph{oi kālo kāpôṛṭā}, \emph{ei nīl jāmāṭā}
  \item the family, respectfully and at a distance — \emph{ei bhāiṭi}, \emph{oi bônṭi}
  \item one and then another — \emph{ekṭā chā, ārekṭā chā}
  \item asking about each, then denying the answer — \emph{eṭā ki?} … \emph{eṭā jôl nôy}
  \item the face — \emph{ei mukhṭā}
  \item the other job of the far head, from chapter twenty-eight — \textbf{\bn{ও}} meaning \emph{also}
\end{itemize}

\begin{tcolorbox}[breakable,colback=teal!4,colframe=teal!35!black,title={Before you move on}]
Which end of the phrase does the pointer sit on, and which end does the classifier sit on? (\textbf{Front} and \textbf{back}.) Can both appear at once? (\textbf{Yes} — and a Bengali speaker usually uses both.) What is \textbf{\bn{এটা}} made of? (\textbf{The near head and the classifier}, with no noun.) And which letter, taught one chapter ago, is inside every one of these? (\textbf{\bn{ট}}.)

Source: \href{https://www.unicode.org/charts/PDF/U0980.pdf}{Unicode Bengali chart}.
\end{tcolorbox}
